\documentclass[10pt, a4paper, landscape]{article}

% -------------------------------------------------
% 宏包引入
% -------------------------------------------------
\usepackage[fontset=mac]{ctex}       % 中文支持
\usepackage{multicol}   % 多分栏
\usepackage{calc}
\usepackage{ifthen}
\usepackage[landscape]{geometry} % 页面设置
\usepackage{amsmath,amsthm,amsfonts,amssymb} % 数学公式
\usepackage{color,graphicx,overpic} % 颜色与图片
\usepackage{hyperref}   % 超链接
\usepackage{enumitem}   % 列表环境控制
\usepackage{titlesec}   % 标题控制
\usepackage{bm}         % 加粗数学符号
\usepackage{xcolor}
\usepackage{tabularx}
\usepackage{array}
\usepackage{booktabs}
\usepackage{color} % 或 xcolor
\usepackage{tikz}       % 绘图
\usetikzlibrary{decorations.pathreplacing, positioning} % 加载brace装饰库
\usetikzlibrary{calc, positioning, arrows.meta, decorations.markings}
\setCJKmainfont{PingFang SC}
\setCJKsansfont{PingFang SC}
\setCJKmonofont{PingFang SC}

% -------------------------------------------------
% 自定义颜色
% -------------------------------------------------

\definecolor{myblue}{HTML}{003153} % 蓝色字
\definecolor{myred}{HTML}{85120F}  % 红色字
\definecolor{mygreen}{HTML}{218254}  % 绿色字
\definecolor{hlblue}{HTML}{ADC1DB} % 蓝色高亮
\definecolor{hlred}{HTML}{740003}  % 红色高亮
\definecolor{hlred40}{HTML}{D5807F}  % 红色高亮40%透明度（在白色背景上模拟）
\definecolor{hlyellow}{HTML}{E3C79F}  % 奢金高亮
\definecolor{hlgreen}{HTML}{BED49D}  % 抹茶绿高亮

% -------------------------------------------------
% 极限空间压缩设置 (核心部分)
% -------------------------------------------------

% 1. 页边距设置为极小 (0.5cm)
\geometry{top=0.5cm,left=0.5cm,right=0.5cm,bottom=0.5cm}

% 2. 去掉段落首行缩进，改为段落间略微留空（可选，这里为了紧凑设为0）
\setlength{\parindent}{0pt}
\setlength{\parskip}{0pt}

% 3. 设置正文基础字体大小为 scriptsize (约8pt)，如果还觉得大，可以改为 \tiny
\renewcommand{\baselinestretch}{0.9} % 压缩行间距
\let\oldfootnotesize\footnotesize
\renewcommand{\footnotesize}{\fontsize{7pt}{8pt}\selectfont}

% 4. 压缩列表环境 (Itemize/Enumerate) 的间距
\setlist{nolistsep} 
\setlist[itemize]{leftmargin=*}
\setlist[enumerate]{leftmargin=*}

% 5. 压缩标题间距
\titleformat{\section}{\bfseries\scriptsize\color{myblue}}{}{0em}{}[\hrule] % 标题带下划线，蓝色，省空间
\titlespacing*{\section}{0pt}{2pt}{1pt} % 上方留2pt，下方留1pt
\titleformat{\subsection}
    [runin] % 不换行
    {\bfseries\scriptsize} % 粗体、scriptsize，黑色字体
    {} % 不显示编号
    {0pt} % 标题与正文间距
    {\hlyellow} % 用hlyellow高亮命令包裹标题
    [] % 标题内容后无内容
\titleformat{\subsubsection}
    [runin] % 不换行
    {\bfseries\tiny} % 粗体、scriptsize，黑色字体
    {} % 不显示编号
    {0pt} % 标题与正文间距
    {\hlgreen} % 用hlgreen高亮命令包裹标题
    [] % 标题内容后无内容
\titlespacing*{\subsection}{0pt}{1pt}{0.5em} % 上方1pt,下方0.5em(水平间距)
\titlespacing*{\subsubsection}{0pt}{1pt}{0.5em} % 上方1pt,下方0.5em(水平间距)

% -------------------------------------------------
% 自定义命令
% -------------------------------------------------
% 颜色字
\newcommand{\red}[1]{\textbf{\textcolor{myred}{#1}}}  
\newcommand{\blue}[1]{\textbf{\textcolor{myblue}{#1}}} 
\newcommand{\green}[1]{\textbf{\textcolor{mygreen}{#1}}}
\newcommand{\entry}[2]{$\bullet$ \textbf{#1}: #2\par\vspace{0.5pt}}
% 高亮
\newcommand{\cbox}[2][yellow]{\begingroup\setlength{\fboxsep}{1pt}\colorbox{#1}{\strut#2}\endgroup}
\newcommand{\hlblue}[1]{\cbox[hlblue]{#1}}
\newcommand{\hlred}[1]{\cbox[hlred40]{#1}}
\newcommand{\hlyellow}[1]{\cbox[hlyellow]{#1}}
\newcommand{\hlgreen}[1]{\cbox[hlgreen]{#1}} 
% 图片插入
\newcommand{\img}[2][0.9\linewidth]{%
    {\par\vspace{1pt}\centering\includegraphics[width=#1]{#2}\par\vspace{1pt}}%
}
% 左图右文 (参数: [图片宽度比例]{图片路径}{右侧文字内容})
\newcommand{\imgleft}[3][0.3]{%
    \noindent\begin{minipage}[t]{#1\linewidth}%
        \vspace{0pt}%
        \includegraphics[width=\linewidth]{#2}%
    \end{minipage}%
    \hfill%
    \begin{minipage}[t]{0.98\linewidth - #1\linewidth}%
        \vspace{0pt}%
        #3%
    \end{minipage}\par\vspace{2pt}%
}
% 左文右图 (参数: [图片宽度比例]{图片路径}{左侧文字内容})
\newcommand{\imgright}[3][0.3]{%
    \noindent\begin{minipage}[t]{0.98\linewidth - #1\linewidth}%
        \vspace{0pt}%
        #3%
    \end{minipage}%
    \hfill%
    \begin{minipage}[t]{#1\linewidth}%
        \vspace{0pt}%
        \includegraphics[width=\linewidth]{#2}%
    \end{minipage}\par\vspace{2pt}%
}
% -------------------------------------------------
% 新增：概念速查表专用命令
% -------------------------------------------------
% 表格容器
\newcommand{\concepttable}[1]{%
    {\setlength{\tabcolsep}{1.5pt}% 局部减小列间距
     \renewcommand{\arraystretch}{0.92}% 局部紧缩行距
     \par\vspace{2pt}{\color{myblue}\hrule height 0.6pt}\vspace{1pt}% 上边框(蓝色,0.6pt粗)
     \noindent\begin{tabular}{@{}p{0.22\linewidth}p{0.76\linewidth}@{}}%
     #1%
     \end{tabular}%
     \vspace{1pt}{\color{myblue}\hrule height 0.6pt}\par\vspace{2pt}}% 下边框(蓝色,0.6pt粗)
}
% 表格行 (参数: {概念名}{解释})
\newcommand{\conrow}[2]{\blue{#1} & #2 \\}


% 定义基于 X 列的自适应等宽列：L (左对齐) 和 C (居中对齐)
\newcolumntype{L}{>{\raggedright\arraybackslash}X}
\newcolumntype{C}{>{\centering\arraybackslash}X}

% ==========================================
% 智能满宽、均匀等分布对比表格环境 \compTable
% 参数说明：
%   [#1]: 可选参数，右侧对比列的格式，默认 C (居中且等宽)
%   {#2}: 必填参数，表格的总列数 (数字)
%   {#3}: 必填参数，表格的完整内容 (包括 header 和 row)
% ==========================================
\newcommand{\compTable}[3][C]{%
    {\setlength{\tabcolsep}{6pt}% 适度保持列间距
     \renewcommand{\arraystretch}{1.15}% 稍微舒展一点的行距
     \par\vspace{3pt}{\color{myblue}\hrule height 0.8pt}\vspace{3pt}% 上边框
     \noindent
     \def\colsnum{#2}%
     % 核心改动：改用 tabularx，宽度锁定为 \linewidth，使用 L 和 C 强制均匀平分列宽
     \begin{tabularx}{\linewidth}{@{} L *{\numexpr\colsnum-1\relax}{#1} @{}}%
         #3%
     \end{tabularx}%
     \vspace{2pt}{\color{myblue}\hrule height 0.8pt}\par\vspace{3pt}}% 下边框
}

% 泛用表格行 \comprow
\newcommand{\comprow}[1]{#1 \\}

% 专门的表头行 \compheader (自带下划线，不使用 \textbf 避免报错)
\newcommand{\compheader}[1]{%
    #1 \\ \noalign{\vspace{1.5pt}{\color{myblue}\hrule height 0.4pt}\vspace{2.5pt}}%
}

% 题干专属方框 (minipage版)
% 用法: \prob{大题号}{小题号}{题干内容}，小题号可空
% 框宽占满行宽，超出行宽时自动换行
\newcommand{\prob}[3]{%
    \par\vspace{3pt}\noindent%
    \fcolorbox{myblue}{myblue!5}{%
        \begin{minipage}{\dimexpr\linewidth-2\fboxsep-2\fboxrule\relax}%
        \vspace{1pt}%
        \textbf{\textcolor{myblue}{[T#1\if\relax\detokenize{#2}\relax\else(#2)\fi]}} #3%
        \vspace{1pt}%
        \end{minipage}%
    }%
    \par\vspace{3pt}%
}


% -------------------------------------------------
% 正文
% -------------------------------------------------
\begin{document}

\tiny

% 三栏布局
\begin{multicols*}{3}

\section{第一章\ 采样定理}

\subsection{前言}

傅里叶变换与傅里叶级数的相关概念复习：

\concepttable{
    \conrow{FT}{$X(j\Omega) = \int_{-\infty}^{\infty} x(t)e^{-j\Omega t}dt$，时域连续非周期 $\mapsto$ 频域连续复函数}
    \conrow{IFT}{$x(t) = \frac{1}{2\pi} \int_{-\infty}^{\infty} X(j\Omega)e^{j\Omega t}d\Omega$，归一化系数 $1/2\pi$}
    \conrow{频率变量}{\textbf{\textcolor{myred}{$\Omega$ 为模拟角频率（rad/s）}}，区别于数字角频率 $\omega$}
}

\concepttable{
    \conrow{FS 展开}{$x(t) = \sum_{m=-\infty}^{\infty} C_m e^{jm\omega_0 t}$，基频 $\omega_0 = \frac{2\pi}{T_0}$}
    \conrow{FS 系数}{$C_m = \frac{1}{T_0} \int_{T_0} x(t)e^{-jm\omega_0 t}dt$}
    \conrow{FS vs.\ FT}{\textbf{\textcolor{myred}{时域周期 $\mapsto$ 频域离散（谱线）；非周期 $\mapsto$ 连续谱}}}
}

\concepttable{
    \conrow{时域卷积}{$x(t) * y(t) \overset{\text{FT}}{\leftrightarrow} X(j\Omega) \cdot Y(j\Omega)$}
    \conrow{系统分析}{$y(t)=x(t)*h(t) \;\Rightarrow\; Y(j\Omega)=X(j\Omega)H(j\Omega)$}
    \conrow{频域卷积}{$x(t) \cdot y(t) \overset{\text{FT}}{\leftrightarrow} \frac{1}{2\pi} X(j\Omega) * Y(j\Omega)$}
    \conrow{采样依据}{理想采样 $x(t) \cdot \delta_T(t)$，频域卷积 $\Rightarrow$ 频谱周期延拓}
}

\textbf{\hlblue{傅里叶变换的充要条件}}

\concepttable{
    \conrow{1. 充分条件（狄利克雷）}{$\int_{-\infty}^{\infty}|x(t)|dt<\infty$，绝对可积}
    \conrow{2. 次强条件（能量有限）}{$\int_{-\infty}^{\infty}|x(t)|^2dt<\infty$}
    \conrow{3. 弱条件（功率有限）}{$\lim_{T\to\infty}\frac{1}{2T}\int_{-T}^{T}|x(t)|^2dt<\infty$}
    \conrow{4. 周期信号}{周期信号为功率有限信号的子类；可用FS求$C_m$，也可用广义FT}
}

\compTable{4}{
    \compheader{$x(t)$ & $X(j\Omega)$ & $x(t)$ & $X(j\Omega)$}
    \comprow{$\delta(t)$ & $1$ & $u(t)$ & $\pi\delta(\Omega)+\frac{1}{j\Omega}$}
    \comprow{$1$ & $2\pi\delta(\Omega)$ & $\operatorname{sgn}(t)$ & $\frac{2}{j\Omega}$}
    \comprow{$e^{j\Omega_0 t}$ & $2\pi\delta(\Omega-\Omega_0)$ & $e^{-at}u(t),\,a>0$ & $\frac{1}{a+j\Omega}$}
    \comprow{$\cos(\Omega_0 t)$ & $\pi[\delta(\Omega-\Omega_0)+\delta(\Omega+\Omega_0)]$ & $\operatorname{rect}\!\left(\frac{t}{\tau}\right)$ & $\tau\operatorname{sinc}\!\left(\frac{\Omega\tau}{2}\right)$}
    \comprow{$\sin(\Omega_0 t)$ & $\frac{\pi}{j}[\delta(\Omega-\Omega_0)-\delta(\Omega+\Omega_0)]$ & $\frac{W}{\pi}\operatorname{sinc}(Wt)$ & $\operatorname{rect}\!\left(\frac{\Omega}{2W}\right)$}
    \comprow{$\sum_{n=-\infty}^{\infty}\delta(t-nT)$ & \multicolumn{3}{L}{$\frac{2\pi}{T}\sum_{k=-\infty}^{\infty}\delta\!\left(\Omega-\frac{2\pi k}{T}\right)$}}
}

\textbf{\hlblue{广义傅里叶变换}}

广义FT将$\delta$函数引入频域，使周期信号（不满足绝对可积条件）也能用FT统一分析。

\concepttable{
    \conrow{核心思想}{\textbf{\textcolor{myred}{引入冲激函数 $\delta(\Omega)$}}，将周期信号的频谱表示为频域冲激串，绕过积分发散问题}
    \conrow{数学基础}{$\mathcal{F}\{e^{j\Omega_0 t}\} = 2\pi\delta(\Omega - \Omega_0)$，即复指数信号的频谱为频域单一冲激}
    \conrow{FT对}{$x(t)=e^{j\Omega_0 t} \;\overset{\text{FT}}{\longleftrightarrow}\; X(j\Omega)=2\pi\delta(\Omega-\Omega_0)$}
}

\concepttable{
    \conrow{周期信号广义FT}{由FS展开 $x(t)=\sum_{m=-\infty}^{\infty}C_m e^{jm\omega_0 t}$，两边取广义FT得：}
    \conrow{频谱表达式}{$X(j\Omega)=\sum_{m=-\infty}^{\infty}2\pi C_m\,\delta(\Omega-m\omega_0)$}
    \conrow{物理含义}{时域周期 $\mapsto$ 频域离散冲激串；谱线位于谐波频率 $m\omega_0$，强度为 $2\pi C_m$}
}


\subsection{模拟信号数字化处理方法}

\subsubsection{采样与量化}

信号依据\textbf{时间轴}与\textbf{幅值轴}的连续/离散属性，分为四类：

\imgleft[0.2]{images/image-2026-06-10-16-20-20.png}{
    \concepttable{
    \conrow{模拟信号}{时间连续、取值连续。物理世界中的基本信号形态}
    \conrow{离散时间信号}{模拟信号经采样（时间离散化）后得到：时间离散、取值连续}
    \conrow{量化阶梯信号}{模拟信号经量化（幅值离散化）后得到：时间连续、取值离散}
    \conrow{数字信号}{离散时间信号再量化（或量化信号再采样），时域与幅值域双重离散：时间离散、取值离散}
}
}

\textbf{\hlblue{信号处理流程}}

$x_a(t)\xrightarrow{\text{预滤波}} \xrightarrow{\text{A/D}} x(n) \xrightarrow{\text{DSP}} y(n) \xrightarrow{\text{D/A}} \xrightarrow{\text{平滑滤波}} y_a(t)$

\concepttable{
    \conrow{预滤波}{限制最高频率，防频谱混叠}
    \conrow{A/D}{采样 + 量化，得数字序列}
    \conrow{DSP}{核心算法，$x(n) \mapsto y(n)$}
    \conrow{D/A}{恢复为阶梯模拟信号}
    \conrow{平滑滤波}{滤除高频镜像，重构 $y_a(t)$}
}

\textbf{\hlblue{时域混叠}}

\concepttable{
    \conrow{连续正弦}{$x_a(t)=A\cos(2\pi F_0 t+\theta)$，$F_0$ 为模拟频率（Hz）}
    \conrow{采样}{设采样频率$F_s=1/T$，令 $t=nT=n/F_s$，得 $x(n)=A\cos(2\pi f_0 n+\theta)$}
    \conrow{数字频率}{$f_0 = \dfrac{F_0}{F_s}$}
    \conrow{混叠}{\textbf{\textcolor{myred}{不同 $F$ 经 $F_s$ 缩放后若对应相同 $f$，则发生频谱混叠}}}
}

\textbf{数学证明}

设 $x_a(t)=A\cos(2\pi F_0 t+\theta)$，另取 $F_k = F_0 + kF_s\ (k \in \mathbb{N})$，则 $x_a'(t)=A\cos(2\pi F_k t+\theta)$。

以 $F_s$ 采样，$t=n/F_s$：
$$x'(n)=A\cos\!\left(2\pi\frac{F_0+kF_s}{F_s}n+\theta\right)=A\cos\!\left(\frac{2\pi nF_0}{F_s}+\theta+2\pi kn\right)$$

由 $\cos(\phi+2\pi kn)=\cos(\phi)$（$k,n\in\mathbb{Z}$），得
$$x'(n)=A\cos\!\left(\frac{2\pi nF_0}{F_s}+\theta\right)=x(n)$$

\concepttable{
    \conrow{结论}{$F_0$ 与 $F_0+kF_s$ 经同一 $F_s$ 采样后所得序列完全相同，即频域混叠在数学上必然发生。}
}

\textbf{\hlblue{数字角频率与模拟角频率}}

\concepttable{
    \conrow{数字角频率 $w$}{$w = \dfrac{2\pi f}{f_s} = 2\pi f_0$}
    \conrow{模拟角频率 $\Omega$}{$\Omega = 2\pi f = w f_s$}
    \conrow{核心映射}{$\displaystyle w = \Omega T_s,\quad \Omega = \frac{w}{T_s}$}
}

和前面的$f$与$f_0$关系类似，分别是模拟信号本身的频率与采样成离散时间信号之后的频率。

\textbf{\hlblue{理想采样（时域建模）}}

\imgleft[0.25]{images/image-2026-06-10-16-53-23.png}{、

\concepttable{
    \conrow{冲激串}{$P_\delta(t) = \sum_{n=-\infty}^{\infty} \delta(t - nT)$}
    \conrow{采样模型}{$x_a(t) \cdot P_\delta(t) = \sum_{n=-\infty}^{\infty} x_a(t)\,\delta(t - nT)$}
    \conrow{采样输出}{$\displaystyle \hat{x}_a(t) = \sum_{n=-\infty}^{\infty} x_a(nT)\,\delta(t - nT)$}
    \conrow{本质}{$\hat{x}_a(t)$ 是强度为 $x_a(nT)$ 的冲激串，\textbf{\textcolor{myred}{仍是连续时间信号}}，是 $x_a(t)$ 与 $x(n)$ 的桥梁}
}
}

\textbf{\hlblue{采样信号频域推导}}

易得采样算子 $P_\delta(t)$ 的频域形式：$\displaystyle P_\delta(j\Omega) = \text{FT}[P_\delta(t)] = \Omega_s \sum_{k=-\infty}^{\infty} \delta(\Omega - k\Omega_s),\quad \Omega_s = \frac{2\pi}{T}$

时域相乘对应频域卷积（由于采样算子的频域也是冲激串，所以卷积结果为复制平移）：

$\hat{X}_a(j\Omega) = \frac{1}{T} \sum_{k=-\infty}^{\infty} X_a(j\Omega) * \delta(\Omega - k\Omega_s) = \frac{1}{T} \sum_{k=-\infty}^{\infty} X_a(j(\Omega - k\Omega_s))$

\entry{核心结论}{\textbf{\textcolor{myred}{采样后频谱是原频谱以 $\Omega_s$ 为周期的周期延拓，幅度缩放 $1/T$}}。$k=0$ 为基带，$k\neq 0$ 为镜像频谱}

\textbf{\hlblue{镜像频谱}}

\imgleft[0.1]{images/image-2026-06-10-17-12-43.png}{

\concepttable{
    \conrow{原频谱}{带限 $X_a(j\Omega)$，截止频率 $\Omega_c$，$[-\Omega_c,\Omega_c]$ 外无分量}
    \conrow{采样序列频谱}{$P_\delta(j\Omega)$ 为间隔 $\Omega_s$ 的离散冲激串}
    \conrow{过采样}{$\Omega_c < \Omega_s/2$，镜像间有隔离带，无混叠，低通滤波（截止 $\Omega_s/2$）可无失真重构}
    \conrow{欠采样}{$\Omega_c > \Omega_s/2$，\textbf{\textcolor{myred}{高频镜像侵入基带，谱线代数叠加}}，频谱结构不可逆破坏}
}

}

\textbf{\hlblue{时域-频域对偶归纳}}

\textbf{\textcolor{myred}{一句话：连续和非周期对应，离散和周期对应。}}

\concepttable{
    \conrow{连续周期}{时域连续 + 周期 $\;\Rightarrow\;$ 频域非周期 + 离散（谱线），工具：FS}
    \conrow{连续非周期}{时域连续 + 非周期 $\;\Rightarrow\;$ 频域非周期 + 连续，工具：FT}
    \conrow{\textbf{\textcolor{myred}{离散周期}}}{时域离散 + 周期 $\;\Rightarrow\;$ 频域周期 + 离散，工具：DFS / DFT}
    \conrow{离散非周期}{时域离散 + 非周期 $\;\Rightarrow\;$ 频域周期 + 连续，工具：DTFT}
}

时域非周期的信号就是包含各种各样的频率，所以频域是连续的。

\entry{注}{离散非周期信号即理想采样后的标准数字信号形态}

\section{第二章\ 时域离散信号及系统频域分析}

\subsection{时域离散信号的傅里叶变换}

\subsubsection{DTFT}

\concepttable{
    \conrow{定义式}{
    \textbf{\textcolor{myred}{$X(e^{j\omega}) = \text{DTFT}[x(n)] = \sum_{n=-\infty}^{\infty} x(n)e^{-j\omega n}$}}}
    \conrow{充分条件（绝对可和）}{$\displaystyle\sum_{n=-\infty}^{\infty} |x(n)| < \infty$，和 Z 变换收敛条件对应}
}

\textbf{\hlblue{IDTFT}}

\concepttable{
    \conrow{定义式}{
    \textbf{\textcolor{myred}{$x(n) = \text{IDTFT}[X(e^{j\omega})] = \frac{1}{2\pi}\int_{-\pi}^{\pi} X(e^{j\omega})e^{j\omega n}d\omega$}}}
}

\textbf{\hlblue{DTFT 的性质}}

主要性质和 Z 变换相同，这里主要列出一些特殊性质。

\compTable{6}{
    \compheader{序列 & 傅里叶变换 & 说明 & 序列 & 傅里叶变换 & 说明}
    \comprow{$nx(n)$ & $j\frac{dX}{d\omega}$ & 频域微分 & $x^*(n)$ & $X^*(e^{-j\omega})$ & 共轭}
    \comprow{$j\operatorname{Im}\{x(n)\}$ & $X_o(e^{j\omega})$ & 虚部 & $\operatorname{Re}\{x(n)\}$ & $X_e(e^{j\omega})$ & 实部}
    \comprow{$x_o(n)$ & $j\operatorname{Im}\{X\}$ & 奇部 & $x_e(n)$ & $\operatorname{Re}\{X\}$ & 偶部}
    \comprow{$x(n)y(n)$ & $\frac{1}{2\pi} \int_{-\pi}^{\pi} X(e^{j\omega})^* Y(e^{-j\omega})$ & & 时域乘积}
}

\subsubsection{离散傅里叶级数}

设 $x(n)$ 为离散周期序列，周期为 $N$。

\textbf{离散傅里叶级数系数；}$\tilde{X}(k)=\sum_{n=0}^{N-1} \tilde{x}(n) e^{-j \frac{2 \pi}{N} k n} \quad 0<k<N-1$

\textbf{DFS 形式：}$\tilde{x}(n)=\operatorname{IDFS}[\tilde{X}(k)]=\frac{1}{N} \sum_{k=0}^{N-1} \tilde{X}(k) e^{j \frac{2 \pi}{N} k n} \quad 0<n<N-1$

\textbf{\hlblue{广义 DTFT}}

$X(e^{j\omega}) = \mathrm{FT}[\tilde{x}(n)] = \frac{2\pi}{N}\sum_{k=-\infty}^{\infty}\tilde{X}(k)\delta\left(\omega-\frac{2\pi}{N}k\right)$

\compTable{4}{%
\compheader{序列 & DTFT & 序列 & DTFT}%
\comprow{$\delta(n)$ & $1$ & $a^n u(n) \quad (|a|<1)$ & $\frac{1}{1 - a e^{-j\omega}}$}%
\comprow{$R_N(n)$ & $\frac{\sin(\omega N/2)}{\sin(\omega/2)} e^{-j\omega(N-1)/2}$ & $u(n)$ & $\frac{1}{1 - e^{-j\omega}} + \displaystyle\sum_{k=-\infty}^{\infty} \pi \delta(\omega - 2\pi k)$}%
\comprow{$x(n)=1$ & $2\pi \displaystyle\sum_{k=-\infty}^{\infty} \delta(\omega - 2\pi k)$ & $e^{j\omega_0 n}$ & $2\pi \displaystyle\sum_{l=-\infty}^{\infty} \delta(\omega - \omega_0 - 2\pi l)$}%
\comprow{$\cos \omega_0 n$ & $\pi \displaystyle\sum_{l=-\infty}^{\infty} [\delta(\omega - \omega_0 - 2\pi l) + \delta(\omega + \omega_0 - 2\pi l)]$ & $\sin \omega_0 n$ & $-j\pi \displaystyle\sum_{l=-\infty}^{\infty} [\delta(\omega - \omega_0 - 2\pi l) - \delta(\omega + \omega_0 - 2\pi l)]$}%
}

\subsubsection{Z 变换}

Z 变换存在的条件是等号右边\textbf{级数收敛，要求级数绝对可和}，即$
\sum_{n=-\infty}^{\infty} \left| x(n) z^{-n} \right| < \infty
$。

\concepttable{%
\conrow{双边 Z 变换}{$X(z) \stackrel{\text{def}}{=} \displaystyle\sum_{n=-\infty}^{\infty} x(n) z^{-n}$}%
\conrow{单边 Z 变换}{$X(z) \stackrel{\text{def}}{=} \displaystyle\sum_{n=0}^{\infty} x(n) z^{-n}$，适用于因果序列（$x(n) = 0, n<0$）}
}

\textbf{\hlblue{Z 变换的收敛域}}

Z 变换的收敛域一般为环状域 $R_{x-} < |z| < R_{x+}$。收敛仅取决于 $|z|$，与相角无关，故\textbf{\textcolor{myred}{收敛域必为以原点为中心的环状区域}}。

\concepttable{%
\conrow{有理分式形式}{$X(z) = \frac{P(z)}{Q(z)}$}%
\conrow{零点}{$P(z_0)=0$ 且 $Q(z_0)\neq 0$，记 $z=z_0$，复平面用 $\circ$ 表示。}%
\conrow{极点}{$Q(z_0)=0$，$|X(z_0)|=\infty$，记 $z=z_0$，复平面用 $\times$ 表示。}%
\conrow{结论}{极点处级数必发散，\textbf{\textcolor{myred}{收敛域内部绝不能包含任何极点}}；边界 $R_{x-}$、$R_{x+}$ 由极点半径决定。}%
}

\textbf{有限长序列收敛域：}

仅在 $n_1\le n\le n_2$ 内非零，$X(z)=\sum_{n=n_1}^{n_2}x(n)z^{-n}$。有限项求和，每项有界则收敛。

\compTable{4}{
    \compheader{条件 & $X(z)$ 幂次特征 & ROC & 排除点}
    \comprow{$n_1<0,\;n_2>0$ & 正、负幂次均有 & $0<|z|<\infty$ & $z=0$、$z=\infty$}
    \comprow{$n_1\ge0$ & 仅负幂次（或常数） & $0<|z|\le\infty$ & $z=0$}
    \comprow{$n_2\le0$ & 仅正幂次（或常数） & $0\le|z|<\infty$ & $z=\infty$}
}

\textbf{右序列收敛域：}

$n \ge n_1$ 非全零，$n < n_1$ 全零。$R_{x-}$ 为最大极点半径。

\compTable{4}{%
    \compheader{条件 & ROC & 排除点 & 说明}%
    \comprow{$n_1<0$（含负时间项） & $R_{x-}<|z|<\infty$ & $z=\infty$ & 含 $z$ 正幂次，$\infty$处发散}%
    \comprow{$n_1=0$（因果序列） & $R_{x-}<|z|\le\infty$ & 无 & 不含 $z$ 正幂次}%
    \comprow{有限长因果序列 & $0<|z|\le\infty$ & $z=0$ & $R_{x-}=0$}%
}

\textbf{左序列收敛域：}

$n \le n_2$ 非全零，$n > n_2$ 全零。$R_{x+}$ 为最小极点半径，与右序列完全对偶。

\compTable{4}{%
    \compheader{条件 & ROC & 排除点 & 说明}%
    \comprow{$n_2>0$（含正时间项） & $0<|z|<R_{x+}$ & $z=0$ & 含 $z$ 负幂次，$0$ 处发散}%
    \comprow{$n_2<0$（特殊左序列） & $0\le|z|<R_{x+}$ & 无 & 仅含 $z$ 正幂次，$0$ 处收敛}%
}

\textbf{双边序列收敛域：}

双边序列可以按线性性质拆分为左序列（反因果）与右序列（因果）：

$X(z) = \sum_{n=-\infty}^{-1} x(n)z^{-n} + \sum_{n=0}^{\infty} x(n)z^{-n}$

\entry{双边 ROC（交集）}{$R_{x-} < |z| < R_{x+}$，即环状区域，\textbf{\textcolor{myred}{若 $R_{x-} \ge R_{x+}$，两区域无交集，$z$ 变换不存在}}。}

\textbf{\hlblue{Z 变换的计算}}

\compTable{6}{
    \compheader{序列 $f(k)$ & $Z$ 变换 $F(z)$ & 序列 $f(k)$ & $Z$ 变换 $F(z)$ & 序列 $f(k)$ & $Z$ 变换 $F(z)$}
    \comprow{$\delta(k)$ & $1$ & $\varepsilon(k)$ & $\dfrac{z}{z-1}$ & $a^k\varepsilon(k)$ & $\dfrac{z}{z-a}$}
    \comprow{$e^{\alpha k}\varepsilon(k)$ & $\dfrac{z}{z-e^{\alpha}}$ & $k\varepsilon(k)$ & $\dfrac{z}{(z-1)^2}$ & $k^2\varepsilon(k)$ & $\dfrac{z(z+1)}{(z-1)^3}$}
    \comprow{$k a^k\varepsilon(k)$ & $\dfrac{az}{(z-a)^2}$ & $\sin(\beta k)\varepsilon(k)$ & $\dfrac{z\sin\beta}{z^2-2z\cos\beta+1}$ & $\cos(\beta k)\varepsilon(k)$ & $\dfrac{z(z-\cos\beta)}{z^2-2z\cos\beta+1}$}
}

\textbf{\hlblue{Z 变换的性质}}

\compTable{6}{
    \compheader{公式 & 公式 & 公式}
    \comprow{$a f_1+b f_2 \leftrightarrow a F_1+b F_2$ & $f(k\pm m) \leftrightarrow z^{\pm m}F(z)$ & $e^{\pm j\omega_0 k}f \leftrightarrow F(e^{\mp j\omega_0}z)$}
    \comprow{$a^k f(k) \leftrightarrow F\!\left(\dfrac{z}{a}\right)$ & $f(-k) \leftrightarrow F(z^{-1})$ & $f_1*f_2 \leftrightarrow F_1F_2$}
    \comprow{$\dfrac{f(k)}{k} \leftrightarrow \displaystyle\int_{z}^{\infty}\dfrac{F(\eta)}{\eta}d\eta$ & $k f(k) \leftrightarrow -z\dfrac{dF(z)}{dz}$ & $\displaystyle\sum_{i=0}^{k}f(i) \leftrightarrow \dfrac{z}{z-1}F(z)$}
    \comprow{$f(0)=\displaystyle\lim_{z\to\infty}F(z)$ & $f(\infty)=\displaystyle\lim_{z\to 1}(z-1)F(z)$& }
}

\textbf{\hlblue{逆Z变换（部分分式展开法）}}

\textbf{整体思路}：把大的有理分式展开为多个简单的分式，每个分式通过查表得到对应的逆变换。

\textbf{左序列处理：}如果分离出来的分式形式为$\frac{bz}{z+a}$，且由收敛域判断其为左序列，则\textbf{\textcolor{myred}{其逆变换为$-b(-a)^{n}u(-n-1)$}}，如果遇到不规则的$u(f(n))$，则为\textbf{\textcolor{myred}{$u(-f(n)-1)$的方式处理为左序列}}。

\subsection{系统函数与频率特性}

DTFT是z变换在单位圆上的特例，因此\textbf{\textcolor{myred}{系统稳定的条件是 $H(z)$ 的所有极点都在单位圆内}}，此时DTFT存在。

\textbf{LTI 离散系统的三种表示方式}：冲激响应，差分方程，系统函数

\concepttable{
    \conrow{因果性}{$n < 0$ 时 $h(n) = 0$，即 H(z)的收敛域表示为 $R_{x^{-}} < |z| \leq \infty$}
    \conrow{稳定性}{$H(z)$ 收敛域包含单位圆（即 DTFT 的存在条件）}
    \conrow{因果稳定性}{$R_{x^{-}} < |z| \leq \infty, \quad 0 < R_{x^{-}} < 1$，即所有极点均在单位圆内}
}

\subsubsection{零极点分布与频率特性}
$H(z)= K_1 z^{N-M} \frac{\prod_{m=1}^M (z - c_m)}{\prod_{k=1}^N (z - d_k)}$，其中$c_m$为零点，$d_k$为极点。

系统稳定时，可化简为：$H(e^{j\omega}) = K e^{j(N-M)\omega} \frac{\prod_{m=1}^M (e^{j\omega} - c_m)}{\prod_{k=1}^N (e^{j\omega} - d_k)} = |H(e^{j\omega})| e^{j \arg[H(e^{j\omega})]}$

\textbf{幅频函数（$|H(e^{j\omega})|$）}：等于增益模长（K），乘以单位圆当前频点 $e^{j\omega}$ 到所有零点的向量长度之积，再除以该频点到所有极点的向量长度之积。

\textbf{注：}

1. \textbf{\textcolor{myred}{位于原点处的零点或极点}}对系统的幅频特性完全没有影响；但它们会在相频函数中引入 $(N-M)\omega$ 的线性相位移。

2. \textbf{\textcolor{myred}{实数系统的所有零点和极点，要么严格位于实轴上，要么必须成对以共轭复数的形式出现}}

3. 若分子多项式 $P(z)$ 与分母多项式 $Q(z)$ 包含完全相同的根，则该零极点在代数上可以相互抵消，从而降低系统的实际阶数

\textbf{\hlblue{区分 IIR 和 FIR}}

\concepttable{
    \conrow{时域判定}{若 $h(n)$ 为\textbf{无穷长序列}，则系统为 IIR 滤波器；
    
    若 $h(n)$ 为\textbf{有限长序列}，则系统为 FIR 滤波器。}
    \conrow{频域判定}{包含至少一个不在原点的极点（非零极点），则系统为 IIR 滤波器；
    
    极点全部位于原点 $z=0$ 处，则系统为 FIR 滤波器。}
}
\subsection{理想采样序列频谱与连续信号频谱的关系}

\section{第三章\ 离散傅里叶变换}

\subsection{离散傅里叶变换的定义}

\subsection{离散傅里叶变换的性质}

\subsection{频域采样}

\subsection{快速傅里叶变换}

\subsection{离散傅里叶变换的应用举例}

\section{第四章\ 时域离散系统的算法结构}

\subsection{时域离散系统的实现}

\subsection{IIR 系统的基本算法结构}

\subsection{FIR 系统的基本算法结构}

\subsection{格形网络}

\subsection{时域离散系统算法结构的计算复杂度}

\section{第五章\ IIR数字滤波器的设计}

\subsection{数字滤波器的基本概念}

\subsection{模拟滤波器设计}

\subsection{IIR 数字滤波器的设计方法}


% --- 手动换栏命令（如果需要强制换列）---
% \columnbreak 

\end{multicols*}

\end{document}